\section{Aula 13}

Nessa aula vimos renormalização de um passo. A ideia é a seguinte,
suponha que dado $R = [10n]\times[n]$ temos
$\Pr(H(R)) > 1 - \eps$. Então a gente fechar um pouco os olhos
e considerar uma percolação em que os sítios são pontos em $n \times \Z^2$
e dizemos que um vertice $v$ está aberto se uma cerca perto de $v$
com retângulos tipo $R$
está aberta. Por exemplo, a cerca do $0$ poderia ser os restângulos
que formam $B_{10n} - B_{9n}$.
Agora, um sítio está aberto nesse modelo com uma probabilidade $1 -
\eps'$ e queríamos dizer
que isso já é maior que os bounds fáceis de percolação de $\Z^2$,
então percolamos.

Esse argumento está quase certo, mas tem um problema enorme, essa
percolação não é independente.
Para isso vamos construir uma teoria sobre percolação dependente.
Precisamos de algumas
definições.

\begin{definition}
	Uma percolação $(Y_x)_{x \in E}$ é dita $k$-independente se para quaisquer dois
	conjuntos de arestas $A$ e $B$ com $d(A,B) > k$ temos $Y_A \perp Y_B$.
\end{definition}

\begin{definition}
	Dadas percolações $Y$ e $X$, dizemos que $Y$ domina estocásticamente $X$
	ou $Y \succ X$ se para toda $f$ crescente
	\[
		\Ex_Y[f] \geq \Ex_X[f].
	\]
\end{definition}

\begin{obs}
	Essa última definição é equivalente a algo bem mais forte. Que
	existe um coupling $Q$ de $Y$ e $X$ tal que $X_e = 1$ implica em $Y_e = 1$.
	Este não é um teorema fácil de mostrar.
\end{obs}

O principal teorema da aula nos diz que se temos um lowerbound para
qualquer aresta estar aberta e supondo que nossa percolação é
$k$-independente para algum
$k$ finito, então ela domina estocásticamente algum produto de Bernoullis.

\begin{theorem}[Liggott-Schonmann-Stacy]
	Para todo $d,k \geq 1$,  existe $\phi : [0,1] \to [0,1]$ com
	$\lim_{p \to 1} \phi(p) = 1$  tal que se $Y$ é uma
	percolação k-independente e $\Pr(Y_x = 1) \geq p$ para todo $x$
	então $Y \succ \text{Ber}(\phi(p))^{\otimes \Z^d}$.
\end{theorem}

\begin{remark}
	Esse teorema só fala algo para $p$ bem perto de $1$. Considere o exemplo em que
	$x \in \Z$, $U_x \sim U[0,1]$ e $Y_x = \id [U_x < U_{x+1}]$. Então é claro
	que $\Pr(Y_x) \geq 1/2$ para todo $x$ e $Y$ é $1$-independente. No entanto,
	$\Pr(X_0 = X_1 = \dots = X_n = 1) = \frac{1}{n!} < p^n$ para todo
	$p$ se $n$ é suficientemente grande. Em particular, $\phi(1/2)$ dada
	pelo teorema deve ser $0$.
\end{remark}

A prova não é nem um pouco intuitiva. Vamos mostrar um lema
específico que vai implicar o teorema.
\begin{lemma}
	Seja $Y$ uma percolação em $\Z^d$ com $\Pr(Y_e = 1) \geq p$.
	Suponha que $r,s \in [0,1]$ são tais que
	\begin{enumerate}
		\item $(1-r)(1-s)^{|B(0,k)|} \geq 1-p$
		\item $(1-r)r^{|B(0,k)|} \geq 1-p$
	\end{enumerate}
	então $YZ^{s} \succ Z^{rs}$ com $Z^{q}$ é uma percolação Bernoulli com
	parâmetro $q$.
\end{lemma}

\begin{remark}
	Para concluir o teorema, note que trivialmente $Y \succ YZ^s \succ
	Z^{rs}$. Além
	do mais, $\phi(p)$ será o maior valor de $rs$ que satisfaz as desigualdades,
	é um pouco chato ver isso, mas é possível fazer a definição
	de forma que ambos vão para $1$ quando $p \to 1$.
\end{remark}

\begin{proof}[Prova do Lema]
	Tá isso aqui é uma esquisitisse e estou seguindo
	\href{https://sci-hub.box/10.1214/aop/1024404279}{o paper original}.
	A ideia é fazer indução e separar no que a gente sabe.
	Dá para checar que pela definição, olhando para só uma aresta $e$, temos que
	\[
		\Pr((YZ^s)_e = 1) \geq ps \geq rs
	\]
	então temos o caso base. Seja $J$ um conjunto finito de arestas de
	tamanho $n$ e supomos
	que provamos o resultado $YZ^s_I \succ Z^{rs}_I$ para $|I| < |J|$. Seja
	\[
		J = \{e_1, e_2, \dots, e_{j}\} \quad \& \quad J' = J \setminus \{e_j\}
	\]
	basta mostrar que dado $\omega \in {0,1}^{|J'|}$
	\[
		\Pr(Y_{e_j} = 1 | (YZ^{s})_{J'} = \omega) \geq r.
	\]
	Se esse for o caso, conseguimos construir um coupling com as
	probabilidades condicionais.
	A ideia aqui é usar o $Z^s$ para esquecermos o valor que $Y$ pode ter tomado.
	Dividimos as arestas de $J'$ em três conjuntos disjuntos,
	\[
		J' = A_0 \cup A_1 \cup M
	\]
	onde $M = \{e_k : d(e_k, e_j) > k\}$ são as arestas independentes.
	$A_0 = B(e_j, k) \cap \{e_k : \omega_k = 0\}$
	são as arestas próximas iguais a $0$ e $A_1$ é o resto, i.e as
	arestas próximas iguais a $1$.
	Pelas condições em $Y$, temos que
	\[
		\Pr(Y_{e_j} = 1 | (YZ^s)_{J'} = \omega) = \Pr(Y_{e_j} = 1 |
			(YZ^s)_{A_0} = \bm{0} \; \land \; (YZ^s)_{A_1} = \bm{1} \; \land
		(YZ^s)_M = \omega_{M}).
	\]
	Pela independência do $Z^s$, isso é o mesmo que
	\[
		\Pr(Y_{e_j} = 1 | (YZ^s)_{A_0} = \bm{0} \; \land \; (Y)_{A_1} =
		\bm{1} \; \land (YZ^s)_M = \omega_{M}).
	\]
	Vamos trabalhar na verdade com o evento oposto,
	$\Pr(Y_{e_j} = 0 | (YZ^s)_{A_0} = \bm{0}, Y_{A_1} = \bm{1}, (YZ^s)_M
	= \omega_M)$ para
	obter nossos bounds. Abrindo a probabilidade condicional, temos
	\[
		\frac{\Pr(Y_{e_j} = 0 \; \cap\; (YZ^s)_{A_0} = \bm{0} \;\cap\;
		Y_{A_1} = \bm{1} \;\cap\; (YZ^s)_M = \omega_M)}
		{\Pr((YZ^s)_{A_0} = \bm{0} \;\cap\; Y_{A_1} = \bm{1} \;\cap\;
		(YZ^s)_M = \omega_M)}
	\]
	Crescendo o evento de cima e diminuindo o de baixo (note
	que vamos esquecer os valores de $Y_{A_1 \cup A_0}$),
	achamos que isso é menor que
	\[
		\frac{\Pr(Y_{e_j} = 0 \;\cap \; (YZ^s)_M = \omega_M)}
		{\Pr(Z^s_{A_0} = \bm{0} \cap Y_{A_1} = \bm{1} \;\cap\; (YZ^s)_M = \omega_M)}
		= \frac{\Pr(Y_{e_j} = 0 | (YZ^s)_M = \omega_M)}
		{\Pr(Z^s_{A_0} = \bm{0} \; \cap \; Y_{A_1} = \bm{1} | (YZ^s)_{M} = \omega_M)}
	\]
	Por hipótese e independência, essa última expressão é igual a
	\[
		\frac{1-p}{(1-s)^{|A_0|} \Pr(Y_{A_1}) = \bm{1} | (YZ^s)_{M} = \omega_M}.
	\]
	E é possível fácilmente cotar essa última igualdade por indução por baixo por
	$r^{|A_1|}$. Portanto, temos que
	\[
		\Pr(Y_{e_j} = 0 | (YZ^s)_{J'} = \omega) \leq
		\frac{1-p}{(1-s)^{|A_0|}r^{|A_1|}}.
	\]
	Por fim, como $(1-r)/(1-p) \geq 1/r^{|B(0,k)|}$ e $(1-r)/(1-p) \geq
	1/(1-s)^{|B(0,k)|}$,
	pela desigualdade geométrica, segue que
	\[
		\frac{1-p}{(1-s)^{|A_0|}r^{|A_1|}} \leq 1-r
	\]
	o que completa a demonstração.
\end{proof}
